\documentclass[11pt,a4paper]{article}
\usepackage[utf8]{inputenc}
\usepackage{amsmath}
\usepackage{amsfonts}
\usepackage{amssymb}
\usepackage{booktabs}
\usepackage{geometry}
\usepackage{hyperref}
\usepackage{listings}
\usepackage{xcolor}

\geometry{top=2.5cm, bottom=2.5cm, left=2.5cm, right=2.5cm}

\definecolor{codeblue}{rgb}{0.1,0.2,0.6}
\definecolor{codegreen}{rgb}{0.1,0.5,0.1}
\definecolor{codegray}{rgb}{0.4,0.4,0.4}

\lstset{
    backgroundcolor=\color{gray!5},
    commentstyle=\color{codegreen},
    keywordstyle=\color{codeblue},
    numberstyle=\tiny\color{codegray},
    stringstyle=\color{purple},
    basicstyle=\ttfamily\small,
    breakatwhitespace=false,         
    breaklines=true,                 
    captionpos=b,                    
    keepspaces=true,                 
    numbers=left,                    
    numbersep=5pt,                  
    showspaces=false,                
    showstringspaces=false,
    showtabs=false,                  
    tabsize=2,
    frame=single
}

\title{Comprehensive Analysis and Technical Blueprint \\ for Modeling Anomalous Heat Conductivity in GOL-3-class Open Traps \\ \large From Microscopic Kinetic Principles to Hybrid MPI+OpenMP Code}
\author{Advanced Plasma Transport Collaboration Workspace}
\date{\today}

\begin{document}

\maketitle

\begin{abstract}
This document serves as an exhaustive consolidation of the physics, computational challenges, debugging paradigms, and production-ready code architectures developed to study the \textbf{anomalous heat conductivity effect} in the GOL-3 multi-mirror plasma trap. Section~\ref{sec:physics} reviews the fundamental microturbulence mechanisms. Section~\ref{sec:open_q} establishes high-impact research fronts and cross-facility applicability. Section~\ref{sec:roadmap} outlines a verified, 6-step hybrid \texttt{MPI}+\texttt{OpenMP} development strategy, detailing foundational physics debug sessions, numerical pitfalls (such as DC space-charge arcs and phase-inversions), and the deployment of a spectral (DFT) Poisson core featuring automated rolling plateau detection diagnostics.
\end{abstract}

\section{The Physical Foundations of GOL-3 Anomalous Transport}\label{sec:physics}
The \textbf{anomalous heat conductivity effect} inside the GOL-3 multi-mirror plasma trap (Budker Institute of Nuclear Physics, Novosibirsk) represents a dramatic \textbf{\(\sim 1,000\)-fold suppression of longitudinal electron thermal conductivity} relative to standard classical predictions. While classical open traps suffer from rapid axial cooling via free-streaming electron thermal transport to the solid exit receivers, GOL-3 acts as a highly effective, self-generated turbulent thermal barrier.

\subsection{Microscopic Mechanism}
The suppression phenomenon unfolds in a multi-stage kinetic sequence during high-power relativistic electron beam (REB) injection (\(\sim 0.8\text{--}1\text{~MeV}\), \(\sim 20\text{--}30\text{~kA}\), \(\sim 8\text{--}12\text{~\(\mu\)s}\)) into a dense deuterium plasma (\(n_e \sim 10^{20}\text{--}10^{22}\text{~m}^{-3}\)):
\begin{itemize}
    \item \textbf{Two-Stream Instability Induction:} The high-current relativistic beam streams through the background thermal core, exciting violent high-frequency electrostatic oscillations known as \textbf{Langmuir microturbulence} (\(\omega \approx \omega_{pe}\)).
    \item \textbf{Enhanced Effective Collision Frequency:} Background plasma electrons scatter intensely off the fields of these localized turbulent wave packets rather than undergoing standard Coulomb collisions with discrete background ions. This inflates the effective collision frequency (\(\nu_{\text{eff}}\)) by orders of magnitude.
    \item \textbf{Choking of Parallel Heat Flux:} Because parallel electron thermal conductivity (\(\kappa_\parallel\)) is inversely proportional to collision frequency, this massive turbulent scattering background chokes off the axial flow of hot tail electrons, yielding an effective suppression factor of \(\kappa_{\text{eff}} \approx 10^{-3} \kappa_{\text{Spitzer}}\).
\end{itemize}

\subsection{Experimental Validation Paradigms}
This microscopic framework has been validated through several high-precision diagnostic channels at BINP:
\begin{itemize}
    \item \textbf{Collective Thomson Scattering (CTS):} Operates at scale thresholds near the Debye length (\(k\lambda_D \le 1\)), directly detecting an immense non-thermal spike in fluctuations at \(\omega \approx \omega_{pe}\) during beam injection, alongside lower-frequency sub-harmonic cascades signifying ion-acoustic coupling.
    \item \textbf{Macroscopic Energy Balances:} Ruby laser Thomson scattering and diamagnetic loop sensors reveal that while the core column stabilizes at multi-keV temperatures (\(T_e \sim 2\text{--}4\text{~keV}\)), it maintains an extremely steep gradient against the cold exit foils. Sustaining such a gradient without instant collapse is macroscopic proof of thermal insulation.
\end{itemize}

\section{Open Research Frontiers and Cross-Facility Applicability}\label{sec:open_q}
Despite decades of experimental data, several critical questions remain open, rendering first-principles kinetic simulation a highly publishable research front:
\begin{itemize}
    \item \textbf{Decoupling Electrostatic and Magnetic Constraints:} Experimentalists trace transport barriers to \textbf{electron density cavities (cavitons)} driven by wave ponderomotive force, which reflect hot electrons. However, self-generated helical magnetic fields (\(\delta \mathbf{B}_\perp\)) also tangle field lines. Dissecting the individual weight of these two mechanisms requires numerical isolation.
    \item \textbf{Return-Current vs. Heat-Flux Asymmetry:} Resolving how cold bulk electrons flow freely to neutralize the beam current while high-velocity tail electrons are strongly scattered and blocked from carrying heat remains a subtle kinetic puzzle.
\end{itemize}
This physics translates directly to other high-interest fusion configurations:
\begin{itemize}
    \item \textbf{Inertial Confinement Fusion (Fast Ignition):} Relativistic beam-core couplings share direct scaling laws with targets at facilities like OMEGA EP and LFEX.
    \item \textbf{Tokamak Runaway Disruptions:} Runaway beams in ITER or SPARC excite matching Langmuir cascades that alter local energy deposition.
    \item \textbf{Modern Mirrors Machine Concepts:} Informing beam-driven or RF-driven turbulence barriers in steady-state machines like GOL-NB, WHAM, or commercial high-field mirror traps.
\end{itemize}

\section{The Incremental MPI+OpenMP Development Blueprint}\label{sec:roadmap}
To isolate these phenomena on locked-down institutional clusters without the installation overhead of exascale libraries, a 6-step standalone \texttt{MPI}+\texttt{OpenMP} roadmap was executed.

\subsection{Development Mapping Summary}
\begin{table}[h]
\centering
\caption{Incremental Implementation Map Summary}
\begin{tabular}{lll}
\toprule
\textbf{Development Step} & \textbf{Targeted Component} & \textbf{Analytical Verification Target} \\
\midrule
Step 1: Serial Physics & Boris-Vay Kinematics & Relativistic Larmor Radius \& Frequency \\
Step 2: Serial Grid & Shape Functions \& Poisson & Electrostatic Plasma Frequency \(\omega_{pe}\) \\
Step 3: OpenMP Threading & Loop Parallelization \& Reductions & Two-Stream Instability Linear Growth Rate \(\Gamma\) \\
Step 4: MPI Scaling & Domain Slicing \& Guard Cells & Particle Migration Stability Across Nodes \\
Step 5: Diagnostics & In-Situ Moments \& Tracers & Anisotropy Thermalization Relaxation Time \\
Step 6: Scale Target & Caviton \& Transport Barriers & Anomalous Suppression Ratio \(\kappa_{\text{eff}} / \kappa_{\text{SH}} \sim 10^{-3}\) \\
\bottomrule
\end{tabular}
\end{table}

\subsection{Computational Resolution Scalings}
\begin{itemize}
    \item \textbf{Spatial Grid Mesh Step (\(\Delta x\)):} Must strictly resolve the electron Debye length to prevent unphysical numerical grid heating. For \(n_e \sim 10^{21}\text{~m}^{-3}\) and initial temperatures \(T_e \sim 10\text{~eV}\), \(\lambda_D \approx 0.74\text{~\(\mu\)m}\), requiring \(\Delta x \le 1.0\,\lambda_D\).
    \item \textbf{Total Domain Boundaries:} Sized to capture several electron skin depths (\(\lambda_s = c/\omega_{pe} \approx 168\text{~\(\mu\)m}\)) to allow full non-linear wave cascades and current filamentation to develop.
    \item \textbf{Grid Sampling Density:} Mandates \(\ge 32\text{--}64\) macro-particles per cell to push the artificial statistical shot-noise collision floor below the true physical plasma turbulence scattering rate.
\end{itemize}

\subsection{Chronology of Core Debugging Sessions}
Transitioning from simple models to an active instability revealed three major numerical pitfalls that required strict physics fixes:
\begin{enumerate}
    \item \textbf{The Direct Current (DC) Space-Charge Blast:} Initial fluid-neutralizing additions caused a dimension/unit mismatch relative to spatial depositions. This triggered a massive unphysical background DC field that acted like a giant battery, forcing particles into boundaries and masking all microscopic waves. \textit{Fix:} Implemented exact dynamic charge neutrality (\(\sum \rho_j = 0\)) and forced a zero-mean electric field condition (\(\sum E_j = 0\)) at every step to shave off the DC ramp.
    \item \textbf{Velocity Phase Eradication:} Applying reflective boundary walls to a unidirectional streaming beam flipped the beam vectors (\(u_x \to -u_x\)). This turned the streaming beam into a counter-streaming matrix, causing rapid Landau damping that wiped out the wave modes. \textit{Fix:} Implemented strict cyclic periodic boundary conditions for both fields and particle tracking coordinates.
\end{enumerate}

\end{document}
